% ============================================================
% cita-vancouver-config.tex
% ------------------------------------------------------------
% SIN IMPLEMENTAR — DELIBERADAMENTE VACÍO
% ------------------------------------------------------------
% vancouver ES EL ÚNICO DE LOS CUATRO ESTILOS DE gbpublisher
% QUE VA POR bibtex Y NO POR biblatex. NO ES UNA DIFERENCIA
% ESTÉTICA: SON DOS MOTORES CON COMANDOS DISTINTOS.
%
% BAJO bibtex NO EXISTEN \parencite, \textcite, \autocite, LA
% VARIANTE ESTRELLADA NI LOS DOS CORCHETES DE PRENOTA Y
% POSTNOTA. SOLO \cite[postnota]{clave}, Y EL PREFIJO, SI LO
% HAY, VA COMO TEXTO LIBRE ANTES DEL COMANDO.
%
% EN LIBROS HAY ADEMÁS UN SEGUNDO PROBLEMA: refsection ES DE
% biblatex, ASÍ QUE LA BIBLIOGRAFÍA POR CAPÍTULO NECESITA OTRO
% MECANISMO. EL CANDIDATO ES chapterbib, QUE SE APOYA EN EL
% .aux POR ARCHIVO QUE GENERA \include, PERO NO ESTÁ
% VERIFICADO Y CAMBIA LA ORQUESTACIÓN DEL COMPILADOR: OBLIGA
% A CORRER bibtex UNA VEZ POR CAPÍTULO, NO UNA POR DOCUMENTO.
%
% MIENTRAS TANTO, docbook-to-latex.xsl CORTA ANTES DE GENERAR
% NADA CUANDO EL ESTILO DEL PROYECTO ES vancouver.
%
% EN REVISTAS LA RAMA EXISTE EN ObtenerPreambuloEmbebido() Y
% NUNCA SE PROBÓ. CUANDO SE IMPLEMENTE ACÁ, VERIFICAR TAMBIÉN
% QUE ExportarBibTeX() EMITA LOS CAMPOS QUE bibtex ENTIENDE:
% journal, year Y address, Y NO LOS DE biblatex —journaltitle,
% date Y location—, QUE bibtex IGNORA EN SILENCIO.
% ============================================================
